% GENERATED FILE. Edit canonical lessons, then run npm run generate:books.
% canonical-source-hash: fnv1a64:e608a6a6482bfd99
% canonical-lessons: SA-C282-aham-prasannah-asmi, SA-C282-bhavitum-icchami, SA-C282-utsukah, SA-R282-feelings-and-hopes

\chapter{How I Feel, and What I Hope For}
\label{ch:sa-how-i-feel-and-what-i-hope-for}
\hlchaptermodality{282}

\begin{chapteropening}
\textbf{By the end of this chapter:} \emph{I can say how I feel, with the adjective agreeing with me (ahaṃ prasannaḥ / prasannā asmi), and what I want to become (ahaṃ vaidyaḥ bhavitum icchāmi; mama svapnaḥ asti …).}

Complete SA-R282-feelings-and-hopes: say how you feel and what you hope to become.
\end{chapteropening}

\section[\texorpdfstring{\sk{अहं} \sk{प्रसन्नः} \sk{अस्मि} (ahaṃ prasannaḥ asmi)}{अहं प्रसन्नः अस्मि (ahaṃ prasannaḥ asmi)}]{\sk{अहं} \sk{प्रसन्नः} \sk{अस्मि} (ahaṃ prasannaḥ asmi) — how I feel}

\label{lesson:SA-C282-aham-prasannah-asmi}



\begin{quote}
Before the new one: say \textquotedblleft{}I am ill\textquotedblright{}. (\emph{ahaṃ rugṇaḥ asmi}.) And a woman? (\emph{ahaṃ rugṇā asmi}.) Feelings work the same way.
\end{quote}

\subsection*{You'll want to know: the feeling follows the speaker}
\begin{quote}
\textbf{\sk{अहं} \sk{प्रसन्नः} \sk{अस्मि।}}
\end{quote}

\begin{quote}
\emph{ahaṃ prasannaḥ asmi.} — \textquotedblleft{}I am pleased.\textquotedblright{}
\end{quote}

\begin{quote}
\textbf{\sk{अहं} \sk{प्रसन्ना} \sk{अस्मि।}}
\end{quote}

\begin{quote}
\emph{ahaṃ prasannā asmi.} — \textquotedblleft{}I am pleased.\textquotedblright{}
\end{quote}

The first is a man speaking, the second a woman: the word for the feeling takes the speaker's gender, \emph{-aḥ} or \emph{-ā}.

\begin{quote}
\textbf{\sk{मीरा} \sk{भीता} \sk{अस्ति।}}
\end{quote}

\begin{quote}
\emph{mīrā bhītā asti.} — \textquotedblleft{}Mira is afraid.\textquotedblright{}
\end{quote}

\begin{quote}
\textbf{\sk{अरुणः} \sk{श्रान्तः} \sk{अस्ति।}}
\end{quote}

\begin{quote}
\emph{aruṇaḥ śrāntaḥ asti.} — \textquotedblleft{}Arun is tired.\textquotedblright{}
\end{quote}

\begin{quote}
\textbf{\sk{ते} \sk{कुपिताः} \sk{सन्ति।}}
\end{quote}

\begin{quote}
\emph{te kupitāḥ santi.} — \textquotedblleft{}They are angry.\textquotedblright{}
\end{quote}

\textbf{\sk{सन्ति}} (\emph{santi}) is \textquotedblleft{}are\textquotedblright{}, for several.

\subsection*{You'll want to know: happy, and what I like}
\textbf{\sk{सुखी}} (\emph{sukhī}), \textquotedblleft{}happy\textquotedblright{}, and \textbf{\sk{दुःखी}} (\emph{duḥkhī}), \textquotedblleft{}unhappy\textquotedblright{}, end in \emph{-ī} for a man.

\begin{quote}
\textbf{\sk{अद्य} \sk{अहं} \sk{सुखी} \sk{अस्मि।}}
\end{quote}

\begin{quote}
\emph{adya ahaṃ sukhī asmi.} — \textquotedblleft{}Today I am happy.\textquotedblright{}
\end{quote}

What you like keeps the frame you know, with \textbf{\sk{मह्यम्}} (\emph{mahyam}), \textquotedblleft{}to me\textquotedblright{}:

\begin{quote}
\textbf{\sk{मह्यं} \sk{संगीतं} \sk{रोचते।}}
\end{quote}

\begin{quote}
\emph{mahyaṃ saṅgītaṃ rocate.} — \textquotedblleft{}I like music.\textquotedblright{}
\end{quote}

A feeling with its reason:

\begin{quote}
\textbf{\sk{अहं} \sk{प्रसन्नः} \sk{अस्मि},}
\end{quote}

\begin{quote}
\emph{ahaṃ prasannaḥ asmi,}
\end{quote}

\begin{quote}
\textbf{\sk{यतः} \sk{अरुणः} \sk{आगतवान्।}}
\end{quote}

\begin{quote}
\emph{yataḥ aruṇaḥ āgatavān.} — \textquotedblleft{}I am pleased, because Arun has come.\textquotedblright{}
\end{quote}

\subsection*{Your turn}
Say these aloud:

\begin{itemize}
\raggedright
  \item \textquotedblleft{}I am afraid\textquotedblright{} — \emph{ahaṃ bhītaḥ asmi}
  \item \textquotedblleft{}I am tired\textquotedblright{} — \emph{ahaṃ śrāntaḥ asmi}
  \item \textquotedblleft{}I am angry\textquotedblright{} — \emph{ahaṃ kupitaḥ asmi}
\end{itemize}

\begin{tcolorbox}[breakable,colback=teal!4,colframe=teal!35!black,title={Before you move on}]
\textquotedblleft{}I am pleased\textquotedblright{}, said by a woman? (\textbf{\sk{अहं} \sk{प्रसन्ना} \sk{अस्मि}}.) \textquotedblleft{}Mira is afraid\textquotedblright{}? (\textbf{\sk{मीरा} \sk{भीता} \sk{अस्ति}}.) \textquotedblleft{}I like music\textquotedblright{}? (\textbf{\sk{मह्यं} \sk{संगीतं} \sk{रोचते}}.)
\end{tcolorbox}
\section[\texorpdfstring{\sk{भवितुम्} \sk{इच्छामि} (bhavitum icchāmi)}{भवितुम् इच्छामि (bhavitum icchāmi)}]{\sk{भवितुम्} \sk{इच्छामि} (bhavitum icchāmi) — what I want to become}

\label{lesson:SA-C282-bhavitum-icchami}



\begin{quote}
Say \textquotedblleft{}I am pleased\textquotedblright{}. (\emph{ahaṃ prasannaḥ asmi}.) Then \textquotedblleft{}I want to go\textquotedblright{}. (\emph{ahaṃ gantum icchāmi}.) Now say what you want to become.
\end{quote}

\subsection*{You'll want to know: \sk{भवितुम्} \sk{इच्छामि}}
\begin{quote}
\textbf{\sk{अहं} \sk{वैद्यः} \sk{भवितुम्} \sk{इच्छामि।}}
\end{quote}

\begin{quote}
\emph{ahaṃ vaidyaḥ bhavitum icchāmi.} — \textquotedblleft{}I want to become a doctor.\textquotedblright{}
\end{quote}

\textbf{\sk{भवितुम्}} (\emph{bhavitum}) is the \emph{-tum} form of \textbf{\sk{भवति}}: \textquotedblleft{}to become\textquotedblright{}. What you become stays in the same case as you: \textbf{\sk{वैद्यः}}, like \textbf{\sk{अहम्}}.

\begin{quote}
\textbf{\sk{सः} \sk{अध्यापकः} \sk{भवितुम्} \sk{इच्छति।}}
\end{quote}

\begin{quote}
\emph{saḥ adhyāpakaḥ bhavitum icchati.} — \textquotedblleft{}He wants to become a teacher.\textquotedblright{}
\end{quote}

\begin{quote}
\textbf{\sk{अहं} \sk{नगरं} \sk{गन्तुम्} \sk{इच्छामि।}}
\end{quote}

\begin{quote}
\emph{ahaṃ nagaraṃ gantum icchāmi.} — \textquotedblleft{}I want to go to the city.\textquotedblright{}
\end{quote}

\subsection*{You'll want to know: \sk{मम} \sk{स्वप्नः} \sk{अस्ति}}
\textbf{\sk{मम} \sk{स्वप्नः} \sk{अस्ति}} (\emph{mama svapnaḥ asti}), \textquotedblleft{}my dream is\textquotedblright{}, opens a hope; \textbf{\sk{यत्}} gives what it is:

\begin{quote}
\textbf{\sk{मम} \sk{स्वप्नः} \sk{अस्ति},}
\end{quote}

\begin{quote}
\emph{mama svapnaḥ asti,}
\end{quote}

\begin{quote}
\textbf{\sk{यत्} \sk{अहं} \sk{वैद्यः} \sk{भविष्यामि।}}
\end{quote}

\begin{quote}
\emph{yat ahaṃ vaidyaḥ bhaviṣyāmi.} — \textquotedblleft{}My dream is that I will become a doctor.\textquotedblright{}
\end{quote}

\begin{quote}
\textbf{\sk{मीरायाः} \sk{स्वप्नः} \sk{अस्ति},}
\end{quote}

\begin{quote}
\emph{mīrāyāḥ svapnaḥ asti,}
\end{quote}

\begin{quote}
\textbf{\sk{यत्} \sk{सा} \sk{नगरं} \sk{गमिष्यति।}}
\end{quote}

\begin{quote}
\emph{yat sā nagaraṃ gamiṣyati.} — \textquotedblleft{}Mira's dream is that she will go to the city.\textquotedblright{}
\end{quote}

\textbf{\sk{मीरायाः}} (\emph{mīrāyāḥ}) is \textquotedblleft{}Mira's\textquotedblright{}.

\subsection*{Your turn}
Say these aloud:

\begin{itemize}
\raggedright
  \item \textquotedblleft{}I want to become a doctor\textquotedblright{} — \emph{ahaṃ vaidyaḥ bhavitum icchāmi}
  \item \textquotedblleft{}I want to become a teacher\textquotedblright{} — \emph{aham adhyāpakaḥ bhavitum icchāmi}
  \item \textquotedblleft{}I want to become a student\textquotedblright{} — \emph{ahaṃ chātraḥ bhavitum icchāmi}
\end{itemize}

\begin{tcolorbox}[breakable,colback=teal!4,colframe=teal!35!black,title={Before you move on}]
\textquotedblleft{}To become\textquotedblright{}? (\textbf{\sk{भवितुम्}}.) \textquotedblleft{}He wants to become a teacher\textquotedblright{}? (\textbf{\sk{सः} \sk{अध्यापकः} \sk{भवितुम्} \sk{इच्छति}}.) \textquotedblleft{}My dream is that …\textquotedblright{}? (\textbf{\sk{मम} \sk{स्वप्नः} \sk{अस्ति}, \sk{यत्} …})
\end{tcolorbox}
\section[\texorpdfstring{\sk{उत्सुकः} (utsukaḥ)}{उत्सुकः (utsukaḥ)}]{\sk{उत्सुकः} (utsukaḥ) — eager}

\label{lesson:SA-C282-utsukah}



\begin{quote}
Before the new one: say \textquotedblleft{}I want to become a doctor\textquotedblright{}. (\emph{ahaṃ vaidyaḥ bhavitum icchāmi}.)
\end{quote}

\subsection*{You'll want to know: \sk{उत्सुकः}}
\textbf{\sk{उत्सुकः}} — \emph{utsukaḥ} — \textquotedblleft{}eager\textquotedblright{}. Like every feeling, it takes the gender of the one who feels it: \textbf{\sk{उत्सुका}} (\emph{utsukā}) for a woman.

\begin{quote}
\textbf{\sk{अहम्} \sk{उत्सुकः} \sk{अस्मि।}}
\end{quote}

\begin{quote}
\emph{aham utsukaḥ asmi.} — \textquotedblleft{}I am eager.\textquotedblright{}
\end{quote}

\begin{quote}
\textbf{\sk{मीरा} \sk{उत्सुका} \sk{अस्ति।}}
\end{quote}

\begin{quote}
\emph{mīrā utsukā asti.} — \textquotedblleft{}Mira is eager.\textquotedblright{}
\end{quote}

\begin{quote}
\textbf{\sk{अहं} \sk{नगरं} \sk{द्रष्टुम्} \sk{उत्सुकः} \sk{अस्मि।}}
\end{quote}

\begin{quote}
\emph{ahaṃ nagaraṃ draṣṭum utsukaḥ asmi.} — \textquotedblleft{}I am eager to see the city.\textquotedblright{}
\end{quote}

\textbf{\sk{द्रष्टुम्}} (\emph{draṣṭum}), \textquotedblleft{}to see\textquotedblright{}, is the \emph{-tum} form of \textbf{\sk{पश्यति}}, from the root √dṛś of \textbf{\sk{द्रक्ष्यति}}.

\begin{cousinweb}
\textbf{\sk{उत्सुकः}} (\emph{utsukaḥ}) is built from \textbf{\sk{उत्}} (\emph{ut}), \textquotedblleft{}apart\textquotedblright{}, and \textbf{\sk{सु}} (\emph{su}), \textquotedblleft{}well\textquotedblright{}, with the ending \emph{-ka}, as Monier-Williams takes it apart; he gives its senses as restless and uneasy, then eager and longing for something.
\end{cousinweb}

\subsection*{Your turn}
Say these aloud:

\begin{itemize}
\raggedright
  \item \emph{utsukaḥ}
  \item \textquotedblleft{}The children are eager\textquotedblright{} — \emph{bālāḥ utsukāḥ santi}
  \item \textquotedblleft{}I am eager to see the city\textquotedblright{} — \emph{ahaṃ nagaraṃ draṣṭum utsukaḥ asmi}
\end{itemize}

\begin{tcolorbox}[breakable,colback=teal!4,colframe=teal!35!black,title={Before you move on}]
What does \textbf{\sk{उत्सुकः}} mean? (\textquotedblleft{}Eager\textquotedblright{}.) \textquotedblleft{}Mira is eager\textquotedblright{}? (\textbf{\sk{मीरा} \sk{उत्सुका} \sk{अस्ति}}.)
\end{tcolorbox}
\section[(review)]{Review — feelings and hopes}

\label{lesson:SA-R282-feelings-and-hopes}



\begin{quote}
Say \textquotedblleft{}my bag is lost\textquotedblright{} and \textquotedblleft{}one ticket, please\textquotedblright{}. (\emph{mama syūtaḥ naṣṭaḥ}, \emph{ekaṃ yānapatraṃ dadātu}.)
\end{quote}

\subsection*{Your turn}
Say these aloud:

\begin{itemize}
\raggedright
  \item I am pleased today, because Mira has come.
  \item Mira is tired, and Arun is afraid.
  \item I like music.
  \item I want to become a teacher.
  \item My dream is that I will see the city. I am eager.
\end{itemize}

\begin{tcolorbox}[breakable,colback=teal!4,colframe=teal!35!black,title={Before you move on}]
\textbf{\sk{श्रान्ता}}: a man or a woman? (A woman.) \textbf{\sk{भवितुम्}} is the \emph{-tum} form of which verb? (\textbf{\sk{भवति}}.) Now the other way round: \textquotedblleft{}I am eager\textquotedblright{}, said by a woman? (\textbf{\sk{अहम्} \sk{उत्सुका} \sk{अस्मि}}.)
\end{tcolorbox}
